% !TeX program = xelatex
\documentclass[10pt,twocolumn]{ctexart}
\usepackage[a4paper,margin=1.7cm,columnsep=0.65cm]{geometry}
\usepackage{amsmath,booktabs}
\usepackage[numbers,sort&compress]{natbib}
\usepackage[colorlinks=true,allcolors=blue]{hyperref}
\newcommand{\tmAP}{\textit{t}-mAP}
\title{Persist4D：长序列感知训练与实体记忆的补充实验}
\author{实验整合稿}
\date{2026年10月11日}
\begin{document}
\maketitle

\section{感知训练与实体记忆的作用范围}

稀疏复访中的分割质量和身份延续依赖不同的信息。
ReScene4D通过联合处理多次扫描学习时间一致的语义实例预测\cite{steiner2026rescene4d}；
持久实体状态则在局部网络之外保存实体描述，使暂时未被观察到的实体仍可参与后续关联。
我们用共同初始化的训练对照检验感知器对长前缀的适应性，再用固定的已训练感知器比较
局部窗口与持久状态。两类实验分别约束感知能力和身份恢复的解释。

实体状态保存特征、类别后验、置信度和生命周期，本文将其称为紧凑实体描述。
固定槽数限制显式状态大小和关联矩阵规模；新扫描的点数仍会影响局部感知耗时。
这一表示不带有统计充分性保证，也不保证在任意长历史或任意场景容量下保留全部信息。

\section{共同初始化与锁定评估}

四条轨迹使用相同的预训练编码器来源和任务参数初始状态。
编码器参数冻结，非编码器任务参数从共同初始化开始训练。
Lightning采用32-true精度配置且关闭TF32；原Concerto的FlashAttention内部仍将QKV转为FP16，
计算后转回原dtype。这一共同保留的模型实现不等同于所有算子均以FP32计算。
双复访基线记为E0；长序列课程基线记为E1；阶段分层注意力采样和查询条件特征反馈
分别记为E2与E3。模块轨迹独立训练，均未从已训练基线继续微调。
四条轨迹都完成29,700次优化器更新，训练随机种子为45。
双卡每卡物理批量为1，累积16个微批次，有效批量为32。
等优化步数并不使双复访与长序列课程具有相同输入点数或计算预算。

检查点由CAL固定，随后才运行SEL及正式确认。
选中的E0、E1、E2和E3分别为5,940、23,760、29,700和2,970步。
较早的选点不表示训练提前终止。
但E3的2,970步检查点及同一步E1控制，完成的更新均处于3RScan仅双复访的课程阶段，
尚未以正概率采样三至五复访训练输入；它们的长输入评价因此仍是该早期检查点的长度外推。
不能将完整轨迹后来接受的长前缀监督归到这个选中检查点上。
模块归因采用同一步数的E1控制，工程比较同时保留最佳CAL选点的E1。
长时汇总量定义为三个总体指标的平均：
\begin{equation}
S_{\mathrm{long}}=\frac{\tmAP_{T=3}+\tmAP_{T=4}+\tmAP_{T=5}}{3}.
\end{equation}
各长度的\tmAP{}由官方评价器对全部指定输入池化计算，不使用场景AP的平均值替代。

正式数据来自3RScan/RIO\cite{wald2019rio}的固定历史验证总体。
PB在$T=2,3,4,5$分别含129个输入，覆盖6个物理参考场景；
LOCAL-T2含154个输入，覆盖46个参考场景；ADDITIONAL在$T=2,3,4$
分别含111、77和32个输入，覆盖40、23和8个参考场景。
所有正式参考场景与本轮TRAIN均无重叠。
PB与LOCAL-T2重叠6个场景，LOCAL-T2与ADDITIONAL重叠40个场景，
因此这些结果不是三个独立外部测试集。一个场景内的顺序和前缀也不是独立样本。

\section{长序列训练和模块对照}

\begin{table*}[t]
\centering
\caption{固定正式总体上的池化任务AP（百分比，越高越好）。所有检查点在正式评估前由CAL锁定；E1的2,970步行是反馈模块的同一步数控制。每行PB为四个长度各129输入，LOCAL-T2为154输入。}
\label{tab:formal-native}
\begin{tabular}{lrrrrrrr}
\toprule
检查点 & 更新步 & PB $T=2$ & PB $T=3$ & PB $T=4$ & PB $T=5$ & PB $S_{\mathrm{long}}$ & LOCAL-T2 \\
\midrule
E0 & 5,940 & 15.6942 & 9.7579 & 6.3207 & 4.3600 & 6.8129 & 21.1993 \\
E1 & 23,760 & 17.7515 & 12.0587 & 7.9845 & 5.9204 & 8.6545 & 27.4173 \\
E1 & 2,970 & 14.2620 & 8.4119 & 6.0703 & 4.9526 & 6.4783 & 20.3069 \\
E3 & 2,970 & 14.7891 & 8.8529 & 6.6723 & 5.3639 & 6.9630 & 20.4407 \\
\bottomrule
\end{tabular}
\end{table*}


长序列训练改善了固定正式总体上的感知结果。
表\ref{tab:formal-native}中，最佳CAL选点的E1相对E0，PB的长时均值提高1.8417个百分点，
四个长度均为正；LOCAL-T2提高6.2180个百分点。
ADDITIONAL的三个长度分别提高5.4719、9.8369和5.6755个百分点。
这些结果支持训练长度对该感知器的影响，不能直接解释持久状态的身份恢复效果。

模块对照体现了不同的收益边界。
在SEL中，阶段采样相对同一步数E1的长时均值提高0.5337个百分点，
但单复访指标下降5.4287个百分点，因而未通过既定的无明显短时退化门槛。
反馈模块相对同一步数E1的SEL长时均值提高0.9005个百分点，进入正式确认。
正式PB中的同一步数差值为0.4847个百分点，LOCAL-T2为0.1338个百分点；
ADDITIONAL的$T=3$则下降0.8474个百分点。
反馈模块增加198,528个参数，这些对照尚不能区分查询条件机制和额外参数的贡献。

最佳基线比较进一步限定了反馈模块的结论。
E3在正式PB的长时均值比E1最佳CAL选点低1.6915个百分点。
同一步数下的正信号与最终工程选点的负差值应同时保留，不能将前者写成对最佳基线的改进。
本轮仅有一个独立训练种子。
按实际训练成本预测，完整的第二种子配对训练需要约1,034.50 GPU小时，
超过1,800 GPU小时总预算的剩余额度，因此未执行独立训练复验。

\section{已训练感知器的记忆对照}

我们固定E0和E1的CAL选中权重，比较完整因果前缀加成对特征--类别关联、
最近双复访窗口加同一关联器，以及最近双复访窗口加固定容量持久状态。
三个部署方式共享每个物理序列的当前扫描；后两个方式共享完全相同的局部候选。
阈值和状态配置在查看新对照结果前固定，未用SEL诊断调整模型。

完整历史方式使用当前因果前缀的联合分割输出计算任务AP，关联器用于跨更新的身份评价。
两个局部方式则积累已提交的各阶段官方候选，使用同一平均置信度规则发布轨迹。
因此，局部持久状态相对局部成对关联的比较隔离了记忆策略；
局部方式相对完整历史的比较还包含感知上下文和轨迹发布方式的差异。
身份率的区间按六个物理参考场景配对重采样，每次对事件计数重新池化计算比率。
场景AP的平均值不用于构造总体AP的区间。
这组桥接按物理序列和因果前缀固定运行顺序及随机数标识，
与正式确认采用的排序协议不同；两张表的AP均来自各自实际预测，不相互替代。

同一物体在完整前缀和局部窗口中的官方类别标签可能不同，
因为原始整理器按所见历史中的首个点确定实例类别，并据此筛选实例。
对129条序列的645个因果前缀检查发现125处窗口依赖差异。
对照统一使用官方因果完整前缀真值，核对当前原始坐标和逐点标注一致，
从局部输入的原始实例编号恢复公共评价掩码，并验证累计真值与完整前缀一致。
真值仅用于评价对齐，不进入网络、候选选择或关联。

两个检查点的129条序列均已完成，每点执行1,032次网络前向；
184份官方充分状态通过哈希核验，40份总体及当前帧指标独立复算一致。
表\ref{tab:trained-memory}显示，在完全共享局部候选的条件下，持久记忆改善了任务AP和gap恢复召回。
E0下的AP提高2.5343个百分点，E1下提高2.2785个百分点；两者仍低于相应完整前缀方式。
完整的$T=2$至$T=5$结果、当前帧指标和事件分母保留于随附CSV。

\begin{table}[t]
\centering
\small
\setlength{\tabcolsep}{3pt}
\caption{固定感知器的五复访记忆对照。AP和恢复召回为百分数。
每行129条序列、6个物理参考场景；恢复计数为正确恢复次数/1,101个gap机会，
包含同场景不同序列的重复事件，不是1,101个独立样本。}
\label{tab:trained-memory}
\begin{tabular}{llrrr}
\toprule
Model & Method & AP & Recall & Correct/gaps\\
\midrule
E0 & FH-B2 & 4.5902 & 4.2688 & 47/1101\\
E0 & LOCAL-B2 & 1.3417 & 4.7230 & 52/1101\\
E0 & LOCAL-B4 & 3.8760 & 20.5268 & 226/1101\\
\midrule
E1 & FH-B2 & 5.9381 & 6.6303 & 73/1101\\
E1 & LOCAL-B2 & 1.1343 & 6.5395 & 72/1101\\
E1 & LOCAL-B4 & 3.4128 & 26.4305 & 291/1101\\
\bottomrule
\end{tabular}
\end{table}

LOCAL-B4相对LOCAL-B2的恢复召回差在E0下为15.8038个百分点，
六场景配对bootstrap的95\%区间为[12.6850,23.0576]；
E1下为19.8910个百分点，区间[16.1261,28.4120]。
采用10,000次参考场景重采样，区间描述该固定面板的事件率不确定性，不能替代独立训练重复。
身份切换率差分别为$-1.3858$和$-0.0470$个百分点，
区间分别为[$-3.6842$,0.1945]和[$-2.2364$,2.4506]，均跨零。
因此不能将恢复召回改善推广为已确认的身份切换率下降。
E1下的局部持久状态AP反而低于E0下的3.8760\%；
完整前缀感知指标的提升未自动转化为固定记忆配置下的任务AP提升。

五复访时，E0的完整前缀与局部持久方式的网络加关联耗时中位数分别为969.30和529.40毫秒，
E1对应为965.14和526.96毫秒。两种权重的完整/局部网络峰值分配显存中位数
分别为3.0910/1.6163 GiB。时间不包含单列的输入搬运和后处理，不能称为端到端总延迟；
两个局部方式复用同一感知前向。
显式持久状态为61,008字节，E0/E1最大占用分别为38/40个槽，总容量100，均无拒绝新实体事件。
所测序列最长五次复访且没有容量饱和，尚不能验证更长历史下的淘汰策略或信息保留能力。


\section{机制与运行检查}

机制面板包含六个CAL五复访序列及其嵌套前缀，覆盖两个物理参考场景。
五种固定模型／干预配置共150次独立进程推理，另有2次仅反向传播的梯度诊断。
全部152项已完成，保存的预测与真值已用于独立重算官方状态。
75份汇总状态分别覆盖25个总体和50个场景--长度组合。

训练后的反馈模块在两个梯度诊断中的梯度范数分别为0.9788与1.4578；
交叉注意力、自注意力、解码FFN和mask head的梯度均有限且非零。
诊断使用原始任务目标与12个辅助输出，不进行优化器更新，编码器没有梯度，所有参数保持不变。
这证明模块连接到实际任务目标，不代表其提供准确率收益。
在28组严格上游对齐的配对中，辅助阶段0--7保持完全一致；
反馈后的四个辅助mask输出及最终mask均改变。
紧接反馈的阶段8类别logits不变，阶段9--11及最终类别logits随后改变，
与先更新mask特征、再影响后续查询更新的插入位置相符。

反馈开启相对关闭在五个长度上的CAL总体AP差分别为
+0.7418, +0.4224, +0.2860, +1.3455, +1.2032个百分点。
这些汇总保留所有原始输入；其中两组存在反馈之前的运行差异，下文单独核验。
在满足严格上游对齐的28组中，当前帧新增和丢失的高IoU对象观测均为9次；
按全部前缀计算分别为5次和6次。
同一物体可能在多个前缀被重复观察，这些计数不是独立物体数。
高IoU诊断按同类候选的最佳IoU是否达到0.5计数，不是官方AP或一对一身份关联。
因此不能把候选变化只叙述为新增成功案例。

阶段采样诊断使用已训练E1的六个五复访输入，在实际层级特征上复现既定采样规则。
72次调用均触发原采样限额，共360个可用阶段位置；均匀和分层规则均未漏掉可用阶段。
分层规则使各阶段采样占比的最大最小差为0，均匀规则的调用均值为0.0612。
这支持配额更均衡，但未在该面板上证明消除了阶段遗漏。
训练轨迹中E2的短时退化仍须保留，不能由采样配额性质替代任务性能。

嵌套前缀记录包含目标不可见、命中、丢失、恢复、类别改变和原始查询槽变化。
其中不可见指目标未出现在当前官方真值中，尚未区分物理缺席与标签筛选；
相邻复访间的原始查询槽改变也不等同于持久跟踪ID切换。


运行顺序测试在同一A40上比较独立进程中的输入A与先处理B再处理A，
其中A为固定CAL双复访输入，B为固定CAL五复访输入，每种顺序重复两次。
E1与E3均使用2,970步权重。
在这个固定输入上，同顺序重复的输入、原始输出、官方预测和AP一致；改变顺序后输入仍一致，输出和预测发生变化。
E1的最大logits差为0.196730，最大mask logits差为1.957217；
E3分别为0.000105和0.001230。这个CAL输入的AP差均为零。
因此保留运行条件限定，单输入的AP不变不能证明整个总体对进程历史不敏感。

更大的机制诊断面板暴露出固定新进程协议的重复性限制。
30组反馈开关配对中，2组在反馈之前的辅助输出就不同，尽管输入与评价真值哈希一致。
我们保留全部30组原始结果，将这2组标为运行差异混杂；其余28组满足严格上游配对条件。
对这两个输入，各模式再独立运行两次，共增加8个前向任务。
原始结果和重复结果构成12个同模式两两比较，其中6个具有上游差异，
但这两个输入的所有同模式官方AP与召回指标相同。
单复访异常输入的同模式最终logits和mask logits最大绝对差分别达到1.5195和19.3663。
这些差异不能被称为逐张量一致，也不能仅因AP相同就判为无影响。

最早保存的辅助阶段已出现mask logits差异；本轮未定位更早的具体算子，
不将原因归于某一个CUDA内核。
重复实验不替换原始配对，也不将偶然对齐的重复结果挑选为新的干预对照。
因此，全30组反馈开关的汇总差值描述实际运行结果，不能单独证明反馈模块的因果收益。
该发现限制了固定协议下的可重复性主张；新增训练须先解决原生数值一致性问题。


部署剖析在同一A40上对每个长度的固定CAL输入预热一次、测量三次。
E1在双复访和五复访上的网络中位耗时分别为438.80和1,929.06毫秒，
峰值分配显存为1.3734和6.1170 GiB。
E3对应为442.21和1,943.03毫秒，显存为1.3741和6.1177 GiB。
其中反馈模块本身的CUDA事件耗时分别为0.509和4.821毫秒。
输入读取、搬运和后处理分别记录；上述值为网络耗时。
五个固定输入不是同一序列的嵌套前缀，不能据此估计普遍的长度增长率。

\section{结论范围与复现}

共同训练结果表明，长序列课程可以改善所测正式总体上的时间一致分割。
阶段采样存在短时退化，反馈模块的同一步数收益依赖评价总体，且低于最佳长序列基线。
同权重记忆对照显示局部AP和恢复召回提高，但尚未追平完整前缀AP，身份切换率差的区间跨零。
更广泛的结论仍需要独立训练种子、等参数训练控制、更多独立复访场景，
以及发生容量饱和和更长历史时的实际评价。

四条训练轨迹的日志版本、固定总体、官方指标充分状态和权重身份随实验产物保留。
五份部署任务权重包含全部非编码器参数与必要缓冲区，使用独立编码器重建，
并已通过实际逐张量重载和远端下载字节核验。
本稿用于整合实测证据，尚非完整CVPR投稿包。

\bibliographystyle{plainnat}
\bibliography{references}
\end{document}
